\chapter{Normalization and ELBO Algebra}

This appendix collects the normalization conditions and algebraic conventions needed to state the evidence decomposition from a fixed joint and a normalized recognition law. Its purpose is to make conditioning sets, Radon--Nikodym requirements, and integrability assumptions available in one proof reference.

\section{Finite-design measure conventions}

For the state-tier results, all latent and observation products are finite and carry the finite product sigma-algebras of Chapter~\ref{ch:probability-semantics}. The configuration space \((\mathsf X,\mathscr X)\) has one fixed design, incidence pattern, and set of identified latent and observation fibers, so every state kernel has the same measurable codomain. Continuous coordinates use the declared Borel base measures, discrete coordinates use counting measures, and mixed coordinates use their product. Probability kernels are normalized before densities are introduced. Whenever a logarithmic density ratio is used, its numerator is absolutely continuous with respect to its denominator and both densities are taken relative to a declared dominating measure. These conventions do not define a measure on a continuum of sections.

\section{Reverse-topological normalization of the directed joint}
\label{sec:directed-normalization-proof}

Fix one design point and a topological order \(v_1,\ldots,v_N\) of the finite DAG in Equation~\eqref{eq:abstract-directed-kernel-factorization}. Under the domination assumptions, write its nonnegative density as
\begin{equation}
f(o,k,m)
=\prod_{\ell=1}^{N}
p_{\theta,v_\ell}^{m}(m_{v_\ell}\mid m_{\operatorname{pa}(v_\ell)},X)
p_{\theta,v_\ell}^{k}(k_{v_\ell}\mid k_{\operatorname{pa}(v_\ell)},m_{v_\ell},X)
p_{\theta,v_\ell}^{o}(o_{v_\ell}\mid k_{v_\ell},m_{v_\ell},X).
\label{eq:appendix-directed-density}
\end{equation}
Tonelli's theorem applies because \(f\geq0\). Integrating all \(o_i\) removes every observation factor, since each is a probability density for every fixed \((k_i,m_i,X)\).

Now traverse the nonroot vertices in reverse topological order. Consider the current vertex \(i\). Every child of \(i\), and hence every descendant of \(i\), occurs later in the forward order and has already been integrated out. After those integrations, the only remaining factor that depends on \(k_i\) as its sampled variable is
\begin{equation}
p_{\theta,i}^{k}(k_i\mid k_{\operatorname{pa}(i)},m_i,X),
\end{equation}
whose integral with respect to \(\nu_i^k(dk_i)\) is one for every fixed conditioning value. Once \(k_i\) is removed, the only remaining factor that contains \(m_i\) as its sampled variable is
\begin{equation}
p_{\theta,i}^{m}(m_i\mid m_{\operatorname{pa}(i)},X),
\end{equation}
whose integral with respect to \(\nu_i^m(dm_i)\) is also one. This pair of integrations removes vertex \(i\) without changing the remaining integrand. The argument covers multiple parents because their coordinates are held fixed while the child kernel is integrated.

After every nonroot has been removed, only roots remain. No edge joins one root as a child of another root, because the child would then have a parent. For each root \(r\), integrate its normalized state bridge \(p_{\theta,r}^{k}(k_r\mid m_r,X)\) and then its normalized model root \(p_{\theta,r}^{m}(m_r\mid X)\). The resulting scalar is one. The kernel statement follows without domination by the same successive-kernel argument. Finally, the declared product over \(a=1,\ldots,M\) multiplies \(M\) unit masses, which proves Equations~\eqref{eq:kernel-total-mass-one} and \eqref{eq:density-normalization}.

\chapter{Gaussian Information and Mean-Field Identities}

This appendix fixes the Gaussian information-form notation and the block partitions used by the restricted recognition controls. It will provide the matrix identities required to compare a correlated Gaussian law with block-product approximations.

\chapter{Gauge Transformations and Verification Oracles}

This appendix specifies the transformation rules and finite-dimensional test constructions used to audit covariance, precision, and probability-ratio calculations. These oracles verify declared formulas under controlled inputs and do not stand in for a claim about the active runtime.

\section{Gaussian kernel gauge oracle}
\label{sec:gaussian-gauge-oracle}

For a nonroot state transition, define its residual
\begin{equation}
\delta_i^k
=k_i-A_i^k\Omega_{ij}^k k_j-B_i m_i-a_i^k.
\label{eq:state-transition-residual}
\end{equation}
Equations~\eqref{eq:generative-link-gauge-laws}--\eqref{eq:generative-bridge-offset-gauge-laws} give \(\delta_i^{k\prime}=g_i^k\delta_i^k\). With \(Q_i^{k\prime}=g_i^kQ_i^k(g_i^k)^\top\), direct inversion gives
\begin{equation}
(Q_i^{k\prime})^{-1}
=(g_i^k)^{-\top}(Q_i^k)^{-1}(g_i^k)^{-1},
\qquad
(\delta_i^{k\prime})^\top(Q_i^{k\prime})^{-1}\delta_i^{k\prime}
=(\delta_i^k)^\top(Q_i^k)^{-1}\delta_i^k.
\label{eq:state-transition-quadratic-invariance}
\end{equation}
The determinant obeys
\begin{equation}
\det Q_i^{k\prime}=\det(g_i^k)^2\det Q_i^k.
\label{eq:state-transition-determinant-law}
\end{equation}
Therefore the transformed Gaussian density acquires the factor \(|\det g_i^k|^{-1}\), which is exactly canceled by the coordinate volume change \(dk_i'=|\det g_i^k|dk_i\). The conditional probability measure, rather than its coordinate density alone, is invariantly defined. Replacing \(k\) by \(m\) proves the model-channel result.

For the observation residual
\begin{equation}
\delta_i^o=o_i-H_i k_i-L_i m_i-d_i,
\label{eq:observation-residual}
\end{equation}
Equation~\eqref{eq:observation-parameter-gauge-laws} gives \(\delta_i^{o\prime}=\delta_i^o\) and \(R_i^{o\prime}=R_i^o\). Hence the observation density itself is unchanged. The same calculations apply to the root kernels after replacing the nonroot means and covariances by Equations~\eqref{eq:gaussian-model-root}--\eqref{eq:gaussian-state-root}. These identities provide a finite matrix oracle: sample invertible represented frame matrices and SPD covariances, transform every parameter by the displayed laws, and test the residual, quadratic-form, determinant, root-moment, and observation-mean equalities to numerical tolerance.

\chapter{Notation}

Table~\ref{tab:magent-notation} consolidates the geometric, probabilistic, and configuration-level types. A symbol in the middle column denotes a domain or mathematical type, not merely an informal role. The table reserves \(q_i\) and \(s_i\) for state- and model-level section values or recognition marginals. If a later random configuration contains measure-valued fields, they are written \(\widehat q_i,\widehat s_i\) and related to state-level marginals by an explicit model or summary map; they are never silently substituted for \(q_i,s_i\).

\small
\begin{longtable}{>{\raggedright\arraybackslash}p{0.17\textwidth}>{\raggedright\arraybackslash}p{0.31\textwidth}>{\raggedright\arraybackslash}p{0.42\textwidth}}
\caption{Typed notation for the geometric, finite-state, and optional configuration constructions.}
\label{tab:magent-notation}\\
\toprule
\textbf{Symbol} & \textbf{Domain or type} & \textbf{Meaning} \\
\midrule
\endfirsthead
\multicolumn{3}{l}{\textit{Table \thetable\ continued}}\\
\toprule
\textbf{Symbol} & \textbf{Domain or type} & \textbf{Meaning} \\
\midrule
\endhead
\midrule
\multicolumn{3}{r}{\textit{Continued on next page}}\\
\endfoot
\bottomrule
\endlastfoot
\(\mathcal C,c,c_a,D\) &
\(\mathcal C\) a smooth base; \(c,c_a\in\mathcal C\); \(D=\{c_a\}_{a=1}^{M}\) &
Contextual base, base points, and finite evaluation design. These are not population-graph vertices.\\
\(V,E,\mathfrak I\) &
\(V\) a finite vertex set, \(E\subseteq V\times V\), \(\mathfrak I\) an interaction complex &
Population incidence and oriented interaction links. The graph or complex is structural data, not the principal-bundle base.\\
\(\mathcal A^i\) &
\(\left(\mathcal U_i;q_i,p_i,s_i,r_i,U_i\right)\), with \(\mathcal U_i\subseteq\mathcal C\) &
Agent object carrying supported local statistical sections and a local group-valued frame coordinate. It is not a vertex label, latent vector, or variational factor alone.\\
\(q_i(c),p_i(c)\) &
\(\mathcal B_{\mathrm{state},c}\subseteq\mathcal P(\mathsf K_{i,c})\) &
State recognition marginal and state prior section values at \(c\).\\
\(s_i(c),r_i(c)\) &
\(\mathcal B_{\mathrm{model},c}\subseteq\mathcal P(\mathsf M_{i,c})\) &
Model recognition marginal and model prior section values at \(c\).\\
\(k_i(c),m_i(c)\) &
\(\mathsf K_{i,c}\) and \(\mathsf M_{i,c}\) &
Latent state and model samples. These are realizations, not probability measures or agents.\\
\(Y(c),Y_D,\mathsf Y_D\) &
\(Y(c)\in\prod_i(\mathsf K_{i,c}\times\mathsf M_{i,c})\) for \(c\in\bigcap_i\mathcal U_i\); \(Y_D\) valued in \(\mathsf Y_D\) &
Population sample at one common-support context, finite-design random element, and its finite product sample space.\\
\(o,\mathsf O_D\) &
\(o\in\mathsf O_D=\prod_{a,i}\mathsf O_{i,c_a}\) &
Finite-design observation and observation sample space.\\
\(X,(\mathsf X,\mathscr X)\) &
\(X\in\mathsf X\); \((\mathsf X,\mathscr X)\) measurable, with fixed \(D\), incidence, and identified fibers &
Structural data conditioned upon by the state model. Only structural values on the fixed measurable design may vary.\\
\(P_\theta,p_\theta\) &
\(P_\theta:(\mathsf X,\mathscr X)\to\mathcal P(\mathsf O_D\times\mathsf Y_D,\mathscr O_D\otimes\mathscr Y_D)\) &
Normalized generative probability kernel and, when dominated, its density defined by \(P_\theta(do,dY\mid X)=p_\theta(o,y\mid X)\nu_D^O(do)\nu_D^Y(dy)\). No factor reads a live recognition marginal.\\
\(Q_X(dY\mid o)\) &
\(Q:(\mathsf O_D\times\mathsf X,\mathscr O_D\otimes\mathscr X)\to\mathcal P(\mathsf Y_D,\mathscr Y_D)\) &
Normalized, generally correlated recognition kernel whose coordinate pushforwards give \(q_i(c_a)\) and \(s_i(c_a)\).\\
\(P_0(dX),R(dX\mid o)\) &
\(P_0\in\mathcal P(\mathsf X,\mathscr X)\); \(R:(\mathsf O_D,\mathscr O_D)\to\mathcal P(\mathsf X,\mathscr X)\), with \(R(\mathord\cdot\mid o)\ll P_0\) where required &
Configuration prior and configuration recognition law in the optional hierarchical tier.\\
\(\rho_0(dX),P_{\mathcal F}(dX)\) &
\(\rho_0,P_{\mathcal F}\in\mathcal P(\mathsf X,\mathscr X)\) when the partition function is finite &
Proper configuration reference law and an optional normalized configuration-Gibbs law.\\
\(\widehat q_i,\widehat s_i\) &
Declared measure-valued coordinates of a random \(X\), if that extension is used &
Configuration-level fields distinguished typographically from the state recognition marginals \(q_i,s_i\); an explicit map must relate the two levels.\\
\(U_i(c),g_i\) &
\(U_i(c),g_i\in G\) &
Local principal-frame coordinate and a principal-frame change. The exponential expression \(U_i=\exp\phi_i\) is only local.\\
\(g_i^k,g_i^m\) &
\(\rho_k(g_i)\in\operatorname{Aut}(\mathsf K)\), \(\rho_m(g_i)\in\operatorname{Aut}(\mathsf M)\) &
State and model actions induced by the same principal group element \(g_i\).\\
\(\Omega_{ij}\) &
\(U_iU_j^{-1}\in G\) on a common local domain &
Regime-I frame comparison; its represented actions are \(\Omega_{ij}^{k}=\rho_k(\Omega_{ij})\) and \(\Omega_{ij}^{m}=\rho_m(\Omega_{ij})\).\\
\(L_{ij}\) &
\(L_{ij}\in G\) on an oriented edge of \(\mathfrak I\) &
Independent Regime-II interaction link with \(L_{ij}'=g_iL_{ij}g_j^{-1}\).\\
\(\Sigma_k,\Sigma_m\) &
\(\mathbb S_{++}^{d_k}\), \(\mathbb S_{++}^{d_m}\) &
Nondegenerate state and model covariances, transforming by represented congruence.\\
\(\Lambda_k,\Lambda_m\) &
\(\Lambda_k=\Sigma_k^{-1}\), \(\Lambda_m=\Sigma_m^{-1}\) &
Local precision tensors, transforming by dual inverse congruence. They are distinct from the later global population precision \(J\).\\
\(J\) &
\(\mathbb S_{++}^{\dim\mathsf Y_D}\) in the finite Gaussian reference regime &
Global correlated population precision used by structured recognition and Gaussian mean-field analysis.\\
\(\nu_D^Y,\nu_D^O\) &
Sigma-finite measures on \(\mathsf Y_D\) and \(\mathsf O_D\) &
Declared finite-product reference measures for densities and Radon--Nikodym ratios.\\
\end{longtable}
\normalsize
